\section{Methodology}
\label{sec:methodology}

We introduce \textsc{StructTransfer}, a framework for predicting and improving the
directed transfer induced by reinforcement learning with verifiable rewards
(RLVR).  Given a base policy $\pi_{\theta_0}$, a candidate source pool
$\mathcal{B}$ with verifiable rewards, and an unlabeled development sample
$\mathcal{A}^{u}$ from a target task, our goal is to estimate whether RLVR on
$\mathcal{B}$ will improve performance on the target and to select a
fixed-budget subset of $\mathcal{B}$ that maximizes this transfer.  The target
answers and rewards are unavailable during representation construction, data
selection, RLVR training, hyperparameter selection, and checkpoint selection;
the held-out target test set is used only for final evaluation.

\textsc{StructTransfer} has four components.  First, it represents each problem
as a \emph{Reasoning Transition Graph} (RTG), which preserves the operations,
dependencies, and control flow of a candidate reasoning process
(Section~\ref{sec:rtg}).  Second, it compares the learnable reasoning supply of
the source with the unresolved reasoning demand of the target through
\emph{Directed Reasoning Coverage} (DRC)
(Section~\ref{sec:drc}).  Third, it uses DRC to predict transfer on previously
unseen source--target pairs (Section~\ref{sec:transfer-prediction}).  Finally,
\emph{DRC-Select} chooses a target-relevant and structurally diverse source
subset without accessing target labels or rewards
(Section~\ref{sec:drc-select}).  The representation and selection procedures
are independent of the policy optimizer and can therefore be combined with
GRPO, DAPO, or other RLVR objectives.

\subsection{Reasoning Transition Graphs}
\label{sec:rtg}

\paragraph{Graph representation.}
Let $x$ denote a problem prompt.  We represent a candidate solution process by
a typed directed graph
\begin{equation}
    G_x = (V_x, E_x, o_x, \rho_x),
\end{equation}
where each node $v\in V_x$ is assigned a reasoning operation
$o_x(v)\in\mathcal{O}$ and each edge $e\in E_x$ is assigned a relation type
$\rho_x(e)\in\mathcal{R}$.  The operation vocabulary $\mathcal{O}$ includes
problem interpretation, information retrieval, variable or state setup,
decomposition, arithmetic computation, symbolic transformation, constraint
propagation, case analysis, search, aggregation, and verification.  The edge
vocabulary distinguishes temporal succession, data dependence, and control
flow.  Consequently, two solutions with similar operation frequencies but
different dependency structures need not receive similar representations.

\paragraph{Label-free extraction.}
For each source or target prompt, we sample $K_{\mathrm{trace}}$ candidate traces
from the same frozen base policy $\pi_{\theta_0}$.  A fixed graph extractor
$\mathcal{E}_{\psi}$ maps the prompt and each visible candidate trace to an RTG:
\begin{equation}
    G_x^{(k)} = \mathcal{E}_{\psi}\!\left(x,\tau_x^{(k)}\right),
    \qquad \tau_x^{(k)}\sim\pi_{\theta_0}(\cdot\mid x).
\end{equation}
The same extraction channel is used for source and target examples so that
representation differences cannot be attributed to using gold solutions on
only one side.  Extraction does not condition on target correctness.  We cache
the resulting graphs and validate a stratified sample against independent
human annotations.  An oracle representation built from verified reference
solutions is used only as a diagnostic upper bound and never by the deployable
selector.

\begin{figure}[t]
    \centering
    \figureorplaceholder{figures/fig2_rtg.pdf}{0.96\columnwidth}{48mm}
    \caption{Illustrative Reasoning Transition Graphs.  Typed nodes encode
    reasoning operations, while typed edges preserve temporal, dependency, and
    control-flow relations that are lost in a bag-of-steps representation.}
    \label{fig:rtg}
\end{figure}

\paragraph{Reasoning motifs.}
We summarize an RTG by a vocabulary $\mathcal{M}$ of typed local subgraphs,
called reasoning motifs.  A motif may encode, for example,
\textsc{Setup}$\rightarrow$\textsc{Transform}$\rightarrow$\textsc{Verify} or a
branching \textsc{CaseSplit} followed by independent derivations and an
\textsc{Aggregate} operation.  The vocabulary and extraction radius are fixed
using meta-training tasks before evaluating a held-out target domain.  Let
$a_{x,m}\in[0,1]$ be the normalized frequency of motif $m$ across the candidate
graphs for $x$.  The motif distribution of a dataset $\mathcal{D}$ is
\begin{equation}
    q_{\mathcal{D}}(m)
    = \frac{1}{|\mathcal{D}|}\sum_{x\in\mathcal{D}} a_{x,m}.
\end{equation}
We measure motif compatibility with a bounded graph kernel
\begin{equation}
    \kappa(n,m)=\exp\!\left(-d_{\mathrm{graph}}(n,m)/\tau\right),
\end{equation}
where $d_{\mathrm{graph}}$ compares node types, typed edges, and structural
attributes.  The temperature $\tau$ is calibrated only on meta-development
tasks.

\subsection{Directed Reasoning Coverage}
\label{sec:drc}

A symmetric similarity score cannot, by itself, explain why transfer from
$\mathcal{B}$ to $\mathcal{A}$ may differ from transfer in the reverse
direction.  DRC instead distinguishes \emph{source supply} from
\emph{target demand} and conditions both on the base policy.

\paragraph{Learnable source supply.}
For a source example $x\in\mathcal{B}$, let $\hat p_x$ be the fraction of
$K_{\mathrm{reward}}$ base-policy rollouts accepted by the source verifier.  We
define its frontier learnability as
\begin{equation}
    \ell_x = 4\hat p_x(1-\hat p_x).
\end{equation}
This quantity is maximal near the empirical decision boundary and vanishes
when all group samples receive the same binary reward, where group-relative
RLVR supplies no within-group advantage signal.  Let
$w_{\mathcal{B}}(x)$ be the prespecified fraction of the fixed RLVR exposure
budget allocated to $x$, with $\sum_x w_{\mathcal{B}}(x)=1$.  The learnable
supply of motif $n$ is
\begin{equation}
    s_{\mathcal{B}}(n)
    = \sum_{x\in\mathcal{B}}
      w_{\mathcal{B}}(x)a_{x,n}\ell_x.
\end{equation}

\paragraph{Unresolved target demand.}
Target rewards are unavailable, so we estimate headroom from disagreement
among base-policy answers.  Let $\widehat P_x(y)$ denote the empirical
distribution of extracted final answers from $K_{\mathrm{target}}$ rollouts,
including a distinguished extraction-failure symbol.  We define
\begin{equation}
    h_x
    = \frac{-\sum_y \widehat P_x(y)\log \widehat P_x(y)}
           {\log K_{\mathrm{target}}},
\end{equation}
with $h_x=0$ when all samples agree.  The unresolved target demand for motif
$m$ is
\begin{equation}
    d_{\mathcal{A}}(m)
    = \frac{1}{|\mathcal{A}^{u}|}
      \sum_{x\in\mathcal{A}^{u}} a_{x,m}h_x.
\end{equation}
This proxy is deliberately label-free.  For analysis only, we also consider an
oracle headroom estimate based on verified base accuracy to quantify the cost
of replacing correctness with disagreement.

\paragraph{Coverage.}
The source coverage of a target motif $m$ is
\begin{equation}
    c_{\mathcal{B}}(m)
    = 1-\exp\!\left(
        -\eta\sum_{n\in\mathcal{M}}
        \kappa(n,m)s_{\mathcal{B}}(n)
      \right),
\end{equation}
where $\eta>0$ controls saturation and is fixed on meta-development tasks.
DRC is the target-demand-weighted coverage
\begin{equation}
    \operatorname{DRC}(\mathcal{B}\!\rightarrow\!\mathcal{A})
    =
    \frac{
      \sum_{m\in\mathcal{M}} d_{\mathcal{A}}(m)c_{\mathcal{B}}(m)
    }{
      \sum_{m\in\mathcal{M}} d_{\mathcal{A}}(m)+\epsilon
    }.
    \label{eq:drc}
\end{equation}
DRC is directional because the source contributes verifier-conditioned
learnability while the target contributes label-free unresolved demand.
Swapping the two datasets changes both quantities even when their marginal
motif frequencies are identical.  It is also model-dependent: the same pair
of tasks may have different transfer potential for base policies with
different reasoning frontiers.

\subsection{Prospective Transfer Prediction}
\label{sec:transfer-prediction}

We calibrate DRC on a collection of meta-training source--target pairs for
which RLVR transfer has already been measured.  Let
$T_{i\rightarrow j}$ be the transfer gain defined in
Section~\ref{sec:problem-formulation}.  Our interpretable predictor is
\begin{equation}
    \widehat T_{i\rightarrow j}
    = \beta_0
      + \beta_1\operatorname{DRC}(\mathcal{B}_i\!\rightarrow\!\mathcal{A}_j)
      + \boldsymbol{\beta}^{\top}\mathbf{z}_{ij}
      + u_{\mathrm{model}}+u_i+u_j,
    \label{eq:transfer-predictor}
\end{equation}
where $\mathbf{z}_{ij}$ contains prespecified controls such as semantic
similarity, prompt-length difference, answer-format compatibility, and
verifier type.  The random effects account for the model, source, and target
when sufficient observations are available; a regularized linear model is
used in smaller regimes.  Neither $T_{i\rightarrow j}$ nor any target outcome
from the held-out domain is used when fitting or selecting the final predictor.
This setup enables prospective leave-one-benchmark and leave-one-domain
evaluation rather than post-hoc correlation on the same transfer matrix.

\subsection{Cross-Benchmark Capability Flow}
\label{sec:capability-flow}

An increase in pass@1 may reflect either broader empirical support or a sharper
distribution over solutions the base policy could already produce.  We
therefore estimate item-level transitions between the base and RLVR policies.
For each target item $x$, we draw $N$ independent completions from policy
$j\in\{0,1\}$, where $j=0$ denotes the base policy and $j=1$ the RLVR policy.
If $c_{j,x}$ completions pass the verifier, we use a Beta posterior
\begin{equation}
    p_{j,x}\mid c_{j,x}
    \sim \operatorname{Beta}\!\left(
        \alpha+c_{j,x},\,\beta+N-c_{j,x}
    \right).
\end{equation}
For a sampling budget $k$, the probability of at least one success is
\begin{equation}
    b_k(p)=1-(1-p)^k.
\end{equation}
Using prespecified confidence and solvability thresholds, an item is classified
as empirically solvable, empirically unsolvable, or uncertain under each
policy.  We define \emph{acquisition} as an unsolvable-to-solvable transition,
\emph{forgetting} as the reverse transition, and \emph{retention} as remaining
solvable.  Among retained items, \emph{sharpening} denotes a posterior-supported
increase in $p_{1,x}-p_{0,x}$.  The net empirical boundary expansion is
\begin{equation}
    \operatorname{NBE}@k
    = \frac{1}{|\mathcal{A}|}
      \sum_{x\in\mathcal{A}}
      \left[
        \mathbb{I}\{U_0(x)\wedge S_1(x)\}
        -\mathbb{I}\{S_0(x)\wedge U_1(x)\}
      \right].
\end{equation}
We report this decomposition together with the full pass@$k$ curve and output
diversity.  We use the term \emph{empirical support expansion}, rather than
absolute capability acquisition, because finite sampling cannot prove that a
base policy assigns exactly zero probability to a correct solution.

\subsection{Target-Label-Free Data Selection}
\label{sec:drc-select}

DRC-Select turns the structural transfer score into a training intervention.  Let
$\mathcal{U}$ be a candidate source pool and let $C$ be a token budget.  Each
candidate $x\in\mathcal{U}$ contributes motif-specific learnable supply
$s_x(n)=a_{x,n}\ell_x$.  For a selected set $S$, we compute target coverage by
replacing $s_{\mathcal{B}}(n)$ in Equation~\ref{eq:drc} with the budget-scaled
accumulated supply
$\sum_{x\in S}[\operatorname{cost}(x)/C]s_x(n)$.  At the full budget this uses
the same exposure normalization as the dataset-level DRC definition.  This
yields a set-valued coverage objective
$F_{\mathrm{cov}}(S;\mathcal{A}^{u})$ with diminishing returns: once a target
motif is well covered, additional near-duplicate source examples provide
little gain.

To prevent the selector from concentrating on a single reasoning template, we
combine target coverage with a target-weighted graph-space facility-location
term.  We first define the structural relevance of candidate $x$ as
\begin{equation}
    r_{\mathcal{A}}(x)
    = \sum_{m,n\in\mathcal{M}}
      d_{\mathcal{A}}(m)\kappa(n,m)a_{x,n}\ell_x.
\end{equation}
The diversity term is then
\begin{equation}
    F_{\mathrm{div}}(S;\mathcal{A}^{u})
    = \frac{1}{\sum_{x\in\mathcal{U}}r_{\mathcal{A}}(x)+\epsilon}
      \sum_{x\in\mathcal{U}}
      r_{\mathcal{A}}(x)
      \max_{y\in S}\operatorname{sim}_{\mathrm{RTG}}(x,y).
\end{equation}
The final objective is
\begin{equation}
    S^* = \arg\max_{S\subseteq\mathcal{U}}
      F_{\mathrm{cov}}(S;\mathcal{A}^{u})
      + \lambda F_{\mathrm{div}}(S;\mathcal{A}^{u})
    \quad\text{s.t.}\quad
      \sum_{x\in S}\operatorname{cost}(x)\leq C,
    \label{eq:drc-select}
\end{equation}
where $\operatorname{cost}(x)$ is the estimated training-token cost.  We use a
cost-aware greedy procedure that repeatedly adds the example with the largest
marginal objective gain per token.  The selected example identifiers are
frozen in an immutable manifest before RLVR begins.  All selection baselines
and DRC-Select subsequently use the same policy optimizer, rollout budget, and
checkpoint rule.  At inference time, the resulting policy requires neither a
retriever nor access to target examples.

\subsection{Computation and Reproducibility}
\label{sec:method-computation}

RTG extraction and motif construction are offline and cached.  With a fixed
motif vocabulary, computing DRC requires a sparse motif-kernel multiplication;
DRC-Select reuses these cached contributions during greedy selection.  The
method changes only the source data distribution and introduces no additional
inference-time computation.  We record the graph-extractor revision, motif
vocabulary, base-policy revision, decoding configuration, dataset hashes, and
selected-example manifest for every run.  All thresholds and calibration
parameters in this section are fixed on meta-development tasks before final
evaluation on a held-out target domain.
